\documentclass{article}\usepackage[T1]{fontenc}\begin{document}
For the split endpoints and common universe sizes of the existing all-length construction, let a be any base arrow and T any ordinary payload natural transformation, not necessarily invertible. TaggedMixed.All.taggedDiagram equates two natural transformations built from the actual SumChange.taggedBase at both payloads and SumChange.taggedPayload at both base objects. The other two faces are the tagged chosen comparison between payload change and base pull, and naturality of ChainPayload.allIterateMap at a. The equation applies across all finite lengths, including zero. Its source is the tagged family of weighted fibres, not the weighted fibre of the single coproduct span. Thus it does not yet prove the mixed diagram between SpanPull.All.evaluationBase and SumChange.evaluationPayload.

At manuscript revision 458ff36a4798e432dee5947afc402f9903a47f14, read lines 929--953 (eq:telescope), then the two operands in SumChange.evaluationPayload and SpanPull.All.evaluationBase. TaggedMixed.Cube.Map permits specified natural transformations for the two payload faces: replacing them with functor equalities would not give the interface needed by the eventual inverse comparison. Its diagram\_iff equates the component law to equality of the two three-face pastes. In Presented.square, the chosen PayloadLift.square is whiskered by the strict fibre presentations and corrected at both ends. At zero, Zero.fullLaw pastes this presentation map with FinalCube.decodeLaw for the identity span; it is not an empty-case-by-reflexivity argument. At successors, leafLaw uses FinalCube.fullDiagram for the native positive evaluation. All.targetApp computes the target naturality arrow of the packed iteration at a given tag; All.taggedLaw inserts the component equation at that tag. All.taggedDiagram retains the actual transformation-level faces. Recheck the presented source comparison and the packed target arrow if either construction changes.

Run the two files under evidence/tagged-mixed-author-fresh using lake env lean. In MarkedFamilies.lean payload change doubles each fibre with a Boolean marker. A marker-flipping endomorphism is not in the image, while the noninvertible base inclusion selects different graded types at its two choices of Boolean point. The final law typechecks on a concrete weighted zero object and an object over tag one. ArrowFaces.lean separately tests the abstract cube against the additive natural numbers: retaining both nonidentity payload faces yields the equality 20 = 20; dropping the left, right, or both gives unequal totals. This tests the allowed abstract interface, not whether the particular tagged components have those numerical values.

For the single-span extension, first prove the mixed law for the actual forward comparison of chosen weighted pulls against compare\_payload, including the componentPull-to-leafPull correction. Transport both base and payload faces to the chosen inverse equivalences with their units and counits; then paste with taggedDiagram. The fixed-payload base unit/composition laws of \#201 do not identify these mixed faces. This tranche does not bundle higher coherences, widen the all-length universe convention, construct flattening or a monad multiplication, or recover strict algebras.
\end{document}
